% ═══════════════════════════════════════════════════════════════════
% ЧАСТЬ V. СТЕНД N=15/30 — ЗАМКНУТАЯ СИСТЕМА
% ═══════════════════════════════════════════════════════════════════
\part{Стенд $N=15/30$: замкнутая система лемм и теорем}

\section{Введение к стенду. Принцип самодостаточности}\label{sec:cyclo-intro}

Настоящая часть посвящена четвёртой ступени лестницы стендов ---
циклотомическим уровням $N=15$ и $N=30$. Первые три ступени пройдены
в частях~I--IV: тор (род $g=1$, четыре спиновые структуры, семейство
$\pi/N$), поверхность K3 (решётка
$\Lambda_{K3}=E_8\oplus E_8\oplus U\oplus U\oplus U$, $b_2=22$,
сигнатура $(3,19)$) и якобиан квартики Клейна (род $g=3$, группа
$\Gamma(2,3,7)$, фазы $\pi/2,\ \pi/3,\ \pi/7$). Каждая ступень
поставляла целочисленную геометрическую тройку: порядок элемента $n$,
целочисленный топологический индекс $B$ и спектральный порог
$\lambda_0$. На ступени $N=15/30$ роль $n$ играет уровень $N$ (порядок
корней единицы), роль $B$ --- целочисленные индексы поляризационных
решёток по блокам проводников (вычисляются в
разделе~\ref{sec:indices}), а спектральная компонента $\lambda_0$ на
данном стенде не привлекается вовсе --- стенд строится без спектральных
гипотез, и это разделение ролей между ступенями является частью
замысла.

Объекты стенда --- кривые Ферма
\begin{equation}\label{eq:fermat}
  C_N \;=\; \{\, (x,y)\in \mathbb{P}^2 \;:\; x^N + y^N = z^N \,\},
  \qquad N = 15,\ 30,
\end{equation}
и их якобианы $\Jac(C_N)$. Роды равны
$g(15)=\tfrac{14\cdot 13}{2}=91$ и $g(30)=\tfrac{29\cdot 28}{2}=406$,
так что период-матрицы стенда имеют размеры $91\times 182$ и
$406\times 812$. Группой симметрий служит $\mu_N\times\mu_N$
(умножения координат корнями из единицы); поле циклотомических
констант --- $\QQ(\zeta_N)$.

\subsection{Принцип самодостаточности}

Изложение подчинено одному требованию: \emph{каждое утверждение либо
доказывается в тексте из первых принципов, либо снабжено явным
вычислительным сертификатом точной арифметики или прямого
интегрирования}. В частности, изложение опирается на собственные
замкнутые конструкции, а не на внешние подгоночные методы:

\begin{keybox}{Замена внешних методов собственными конструкциями}
\begin{center}
\small
\begin{tabularx}{\textwidth}{@{}>{\raggedright\arraybackslash}p{0.44\textwidth} >{\raggedright\arraybackslash}X@{}}
\toprule
\textbf{Метод раннего плана} & \textbf{Замена в настоящем тексте} \\
\midrule
«Общая фаза» через L-значения характеров Гекке &
Фаза не ищется: она явно входит в замкнутую форму
$P=\tfrac1N\zeta_N^{ra+sb}\Omega_{a,b}$ (лемма~\ref{lem:closedform}) \\
PSLQ на $N=30$ &
Соотношения выводятся из $\Ga$-тождеств
(леммы~\ref{lem:reflection},~\ref{lem:gauss},~\ref{lem:relations}) \\
Численная проверка соотношений Римана на $91\times182$, $406\times812$ &
Теорема~\ref{th:riemann}: тождественное доказательство; числа ---
одноразовая санити-строка вне текста \\
Сборка симплектического базиса численно &
Комбинаторный базис решётчатого графа накрытия
(леммы~\ref{lem:genus},~\ref{lem:sympl}) \\
Апелляция к самосопряжённому оператору &
Используются только решётка циклов с формой пересечений и интегралы
форм (раздел~\ref{sec:nooperator}) \\
\bottomrule
\end{tabularx}
\end{center}
\end{keybox}

Раздел устроен как замкнутая цепочка: комбинаторика циклов
(раздел~\ref{sec:cycles}) $\to$ замкнутые $\Ga$-формы периодов
(раздел~\ref{sec:periods}) $\to$ $\Ga$-тождества
(раздел~\ref{sec:gammaid}) $\to$ соотношения Римана как тождество
(раздел~\ref{sec:riemann}) $\to$ $\Ga$-нормировка и дробные доли
$\pi$ (раздел~\ref{sec:normalization}) $\to$ поляризация и индексы
решётки Ходжа (раздел~\ref{sec:indices}) $\to$ протокол и статусы
(раздел~\ref{sec:protocol}). Каждая лемма используется в последующих
теоремах; никаких внешних «остатков» цепочка не содержит.

\section{Комбинаторика циклов (лемма~0)}\label{sec:cycles}

В этом разделе фиксируются два топологических объекта, на которых
держится весь стенд: решётка циклов $H_1(C_N,\ZZ)$ и форма пересечений
$J$ на ней. Принципиально, что для их определения не требуется ничего,
кроме комбинаторики накрытия: ни аналитики, ни спектральной теории,
ни каких-либо операторов. Аналитика подключается только в
разделе~\ref{sec:periods} и садится на уже готовую решётку.

\begin{lemma}[род кривой Ферма]\label{lem:genus}
Пусть $N\ge 4$ и $C_N=\{x^N+y^N=z^N\}\subset\mathbb{P}^2$. Проекция
$\pi_x\colon C_N\to\mathbb{P}^1$ на координату $x$ есть морфизм степени
$N$, неразветвлённый над $\infty$ и разветвлённый ровно над $N$ точками
$\zeta_N^k$ ($k=0,\dots,N-1$), причём индекс ветвления в каждой точке
$P_k=(\zeta_N^k,0)$ равен $N$. Род кривой равен
\[
  g \;=\; \frac{(N-1)(N-2)}{2};
  \qquad
  g(15)=91,\quad g(30)=406 .
\]
\end{lemma}

\begin{proof}
Над произвольной точкой $x_0\in\mathbb{P}^1$ уравнение
$y^N=1-x_0^N$ имеет ровно $N$ решений с учётом кратности, поскольку
многочлен $y\mapsto y^N-(1-x_0^N)$ имеет степень $N$; значит,
$\deg\pi_x=N$. В точке $P_k=(\zeta_N^k,0)$ имеем
$1-x^N=-N\zeta_N^{k(N-1)}(x-\zeta_N^k)+\dots$, то есть локально
$y^N=c\,(x-\zeta_N^k)$ с $c\ne0$: параметром служит $y$, а
$x-\zeta_N^k\sim c^{-1}y^N$, поэтому индекс ветвления равен $N$.
Над $\infty$ ветвления нет: в однородных координатах при $z=0$
уравнение $X^N+Y^N=0$ имеет $N$ различных решений, и в каждом из них
проекция неразветвлена (локальный параметр $v=z/Y$ регулярного ранга,
см.~\eqref{eq:infinity}). Формула Римана--Гурвица даёт
\[
  2g-2 \;=\; N\,(2\cdot 0-2) \;+\; \sum_{p\in C_N}(e_p-1)
  \;=\; -2N \;+\; N\,(N-1) \;=\; N^2-3N,
\]
откуда $2g=(N-1)(N-2)$ и $g=(N-1)(N-2)/2$. Подстановка $N=15$ и
$N=30$ даёт $91$ и $406$.
\end{proof}

\begin{lemma}[решётчатый граф накрытия]\label{lem:gridgraph}
Положим $t=x^N$ и рассмотрим отображение
$\varphi=(x^N,y^N)\colon C_N\to\mathbb{P}^1$; над дополнением
$\mathbb{P}^1\setminus\{0,1,\infty\}$ это неразветвлённое накрытие
степени $N^2$ с группой листов $\mu_N\times\mu_N$ (лист --- пара
ветвей $x=\zeta_N^r t^{1/N}$, $y=\zeta_N^s(1-t)^{1/N}$). Пусть
$C_N^\circ=C_N\setminus\varphi^{-1}\{0,1,\infty\}$; тогда $C_N^\circ$
деформационно ретрактируется на граф $\Gamma_N$ --- решётчатый граф
$\ZZ/N\times\ZZ/N$ с $N^2$ вершинами и $2N^2$ рёбрами (два остовных
направления, отвечающие двум образующим $\sigma_0,\sigma_1$ группы
$\pi_1(\mathbb{P}^1\setminus\{0,1,\infty\})\cong F_2$).
Цикломатическое число графа равно $N^2+1$.
\end{lemma}

\begin{proof}
Отождествим $\mathbb{P}^1\setminus\{0,1,\infty\}$ с цилиндром с двумя
разрезами; его фундаментальная группа свободна с образующими
$\sigma_0$ (обход вокруг $0$) и $\sigma_1$ (обход вокруг $1$).
Выберем два пути $\ell_0,\ell_1$ на базе от опорной точки $t_0$ к
$t_1$, реализующие $\sigma_0$ и $\sigma_1$. Каждому листу
$(r,s)\in\mu_N\times\mu_N$ сопоставим подъёмы $\ell_0^{(r,s)}$ и
$\ell_1^{(r,s)}$ --- это даёт $N^2$ вершин (концы подъёмов в слое над
$t_0$) и $2N^2$ рёбер. Монодромия листов проста: обход $\sigma_0$
переводит ветвь $t^{1/N}\mapsto\zeta_N t^{1/N}$, то есть
$(r,s)\mapsto(r+1,s)$; обход $\sigma_1$ действует как
$(r,s)\mapsto(r,s+1)$. Тем самым граф $\Gamma_N$ --- в точности
координатная решётка на торе $(\ZZ/N)^2$. Накрытие над дополнением
трёх точек неразветвлено, поэтому эйлерова характеристика
$\chi(C_N^\circ)=N^2\cdot\chi(\mathbb{P}^1\setminus\{0,1,\infty\})
=N^2(-1)$, что совпадает с подсчётом по графу:
$V-E=N^2-2N^2=-N^2$. Цикломатическое число связного графа равно
$E-V+1=N^2+1$.
\end{proof}

\begin{remark}
Граф $\Gamma_N$ --- это в точности «скелет» в смысле бинарного
кодирования геометрии (раздел~\ref{sec:binary}): вершины --- биты
(положения), рёбра --- переходы, замкнутые пути по рёбрам --- циклы.
Никакого дополнительного содержания в «симплектический базис» не
вкладывается: это чтение нашего же скелета в стандартной
алгебраической форме.
\end{remark}

\begin{lemma}[замыкание циклов и точная последовательность]\label{lem:capping}
Ранги целочисленных гомологий связаны точной последовательностью
\[
  0\;\longrightarrow\; K\;\longrightarrow\;
  H_1(C_N^\circ,\ZZ)\;\longrightarrow\;H_1(C_N,\ZZ)\;\longrightarrow\;0,
\]
где $\rank H_1(C_N^\circ,\ZZ)=N^2+1$ (цикломатическое число графа
$\Gamma_N$), $\rank H_1(C_N,\ZZ)=2g=(N-1)(N-2)$, а ядро $K$ порождено
$3N$ петлями вокруг проколотых точек при единственном соотношении
$\sum(\text{петли})=0$, так что $\rank K=3N-1$.
\end{lemma}

\begin{proof}
Для кривой с $m=3N$ проколами стандартная гомологическая
последовательность пары $(C_N,C_N^\circ)$ даёт
$H_1(C_N^\circ)\to H_1(C_N)\to\widetilde H_0(C_N\setminus C_N^\circ)
\to 0$, где правый член --- $\ZZ^{3N}/\ZZ$ ранга $3N-1$. Ранг левого
члена равен цикломатическому числу $N^2+1$ (лемма~\ref{lem:gridgraph}),
ранг $H_1(C_N,\ZZ)$ равен $2g=(N-1)(N-2)$ (лемма~\ref{lem:genus}).
Согласование рангов: $(N^2+1)-(N^2-3N+2)=3N-1$ --- совпадает с рангом
правого члена, поэтому последовательность точна по построению.
Описание ядра: малая петля вокруг точки слоя над $0$ (то есть точки
$x=0$, $y^N=1$) есть замыкание подъёма $\sigma_0$-петли, пройденной
$N$ раз (после $N$ обходов вокруг $t=0$ ветвь $t^{1/N}$ возвращается
на исходный лист); аналогично для $t=1$ и $t=\infty$.
\end{proof}

\begin{lemma}[симплектический базис]\label{lem:sympl}
На $H_1(C_N,\ZZ)$ форма пересечений $J(\gamma,\gamma')$ ---
кососимметричная унимодулярная билинейная форма. Существует базис
\[
  \mathfrak{B}=\{a_1,\dots,a_g,\; b_1,\dots,b_g\}
\]
такой, что $J(a_i,a_j)=J(b_i,b_j)=0$ и $J(a_i,b_j)=\delta_{ij}$, то
есть матрица формы в нём есть
$J=\begin{pmatrix}0&I_g\\-I_g&0\end{pmatrix}$.
\end{lemma}

\begin{proof}
$C_N$ --- замкнутое ориентируемое многообразие вещественной размерности
два, поэтому по двойственности Пуанкаре $J$ невырождена и
унимодулярна на $H_1(C_N,\ZZ)\cong\ZZ^{2g}$. Классификация замкнутых
ориентируемых поверхностей даёт фундаментальный $4g$-угольник с
граничным словом
$a_1b_1a_1^{-1}b_1^{-1}\cdots a_gb_ga_g^{-1}b_g^{-1}$; рёбра этого
многоугольника, отождествлённые попарно, и дают искомые циклы, а
чтение пересечений рёбер на границе даёт матрицу $J$ (пересечение
$a_i$ с $b_i$ равно $+1$, все прочие пары --- $0$). Существование
фундаментального многоугольника --- стандартная теорема о
классификации поверхностей, опирающаяся только на комбинаторику
триангуляций; никакая аналитика в доказательстве не участвует.
\end{proof}

\begin{statusbox}[статус раздела комбинаторики]
Леммы~\ref{lem:genus}--\ref{lem:sympl} доказаны полностью; опираются
они только на формулу Римана--Гурвица, эйлерову характеристику,
гомологическую последовательность пары и классификацию поверхностей ---
весь инструментарий топологический. Тождественная перепись листов,
рёбер и граней для $N=15$ и $N=30$ (суммы рангов, согласование
$N^2+1=(N-1)(N-2)+(3N-1)$) проверена точной целочисленной арифметикой
в скрипте-протоколе (проверка V1, раздел~\ref{sec:protocol}); никакой
численной сборки базиса не производится и не требуется.
\end{statusbox}

\section{Замкнутые формы периодов (лемма~1)}\label{sec:periods}

\subsection{Действие $\mu_N\times\mu_N$ и собственные формы}

Группа $\mu_N\times\mu_N$ действует на $C_N$ умножением координат:
$(\xi,\eta)\cdot(x,y)=(\xi x,\eta y)$. Это действие определено над
$\ZZ$ и коммутирует со всем дальнейшим; его собственные пространства
организуют и формы, и циклы.

\begin{lemma}[базис голоморфных форм]\label{lem:forms}
Формы
\[
  \eta_{i,j} \;=\; \frac{x^i\,y^j\,dx}{y^{\,N-1}},
  \qquad i,j\ge 0,\quad i+j\le N-3,
\]
голоморфны на $C_N$; их число равно $(N-1)(N-2)/2=g$, и они образуют
базис пространства голоморфных дифференциалов $H^0(C_N,\Omega^1)$.
Каждая $\eta_{i,j}$ --- собственная форма действия $\mu_N\times\mu_N$
с характером $(a,b)=(i+1,\,j+1)$:
\begin{equation}\label{eq:eigen}
  (\xi,\eta)^*\,\eta_{i,j} \;=\; \xi^{\,i+1}\eta^{\,j+1}\,\eta_{i,j}.
\end{equation}
\end{lemma}

\begin{proof}
\emph{Голоморфность.} В точке $P_\zeta=(\zeta,0)$, где $\zeta^N=1$,
из $x^N+y^N=1$ следует $x=\zeta(1-y^N)^{1/N}=\zeta(1-\tfrac{y^N}{N}+
\dots)$, откуда $dx\sim-\zeta\,y^{N-1}dy$, и значит
$\eta_{i,j}\sim-\zeta^{\,i+1}y^{\,j}\,dy$ --- дифференциал регулярного
ранга. В точках бесконечности перейдём к координатам $u=X/Y$,
$v=Z/Y$ (тогда $u^N+1=v^N$, локальный параметр $v$, и $u$ голоморфна
по $v$); поскольку $x=u/v$, $y=1/v$, прямой подсчёт даёт
\begin{equation}\label{eq:infinity}
  \eta_{i,j} \;=\; u^{\,i}\,(v\,du-u\,dv)\,v^{\,N-3-i-j},
\end{equation}
и форма регулярна в точности при $i+j\le N-3$.
\emph{Число.} Число пар $(i,j)$ с $i,j\ge0$, $i+j\le N-3$ равно
$\sum_{k=0}^{N-3}(k+1)=(N-2)(N-1)/2=g$ по лемме~\ref{lem:genus}.
\emph{Независимость.} Формы $\eta_{i,j}$ с различными парами
$(i+1,j+1)\in(\ZZ/N)^2$ имеют различные характеры действия
по~\eqref{eq:eigen}. Собственные векторы различных характеров конечной
абелевой группы линейно независимы: если
$\sum_{(i,j)}c_{i,j}\eta_{i,j}=0$ --- нетривиальная зависимость, то,
применив проекционный оператор на характер $(a,b)$ --- среднее по
группе с весом $\zeta_N^{-(ua+vb)}$ --- получим
$c_{a-1,b-1}\,\eta_{a-1,b-1}=0$, противоречие. Итого $g$ независимых
голоморфных форм --- базис.
\end{proof}

\begin{definition}[индекс набора и проводник]\label{def:conductor}
Для пары $(a,b)\in T_N$, где
$T_N=\{(a,b)\in\ZZ^2:\ 1\le a,\ b,\ a+b\le N-1\}$, положим
\[
  g_0=\gcd(N,a,b),\qquad d=\frac{N}{g_0}.
\]
Число $d$ называется \emph{проводником} характера $(a,b)$. Множество
$T_N$ находится в биекции с базисом форм леммы~\ref{lem:forms}
посредством $(a,b)=(i+1,j+1)$; при этом
$|T_N|=(N-1)(N-2)/2=g$.
\end{definition}

Перепись характеров по проводникам (точные целочисленные таблицы
$h_d=\#\{(a,b)\in T_N:\ d(a,b)=d\}$ и вывод через функцию Мёбиуса)
приведена в приложении~\ref{app:census}; итог:
\[
  N=15:\quad h_3=1,\ h_5=6,\ h_{15}=84,\quad \textstyle\sum h_d=91;
\]
\[
  N=30:\quad h_3=1,\ h_5=6,\ h_6=9,\ h_{10}=30,\ h_{15}=84,\
  h_{30}=276,\quad \textstyle\sum h_d=406.
\]

\subsection{Сведение периодов к бета-интегралу}

\begin{lemma}[замкнутая форма периода]\label{lem:closedform}
Пусть $(a,b)\in T_N$ и $\eta=\eta_{a-1,b-1}$ --- соответствующая
собственная форма. В координате $t=x^N$ справедливо тождество
\begin{equation}\label{eq:beta}
  \eta_{i,j} \;=\; \frac{1}{N}\;
  t^{\frac{a}{N}-1}\,(1-t)^{\frac{b}{N}-1}\,dt,
  \qquad a=i+1,\ b=j+1 .
\end{equation}
Для подъёма $\gamma$ пути $\delta\subset\mathbb{P}^1\setminus
\{0,1,\infty\}$ с данными обхода $(r,s)\in(\ZZ/N)^2$ период имеет
замкнутую форму
\begin{equation}\label{eq:closedform}
  P_{a,b}(r,s) \;=\; \int_\gamma \eta
  \;=\; \frac{1}{N}\,\zeta_N^{\,ra+sb}\;\Omega_{a,b},
\end{equation}
где
\begin{equation}\label{eq:omegaab}
  \Omega_{a,b}\;=\;\Bfun\!\Bigl(\frac{a}{N},\frac{b}{N}\Bigr)
  \;=\;\frac{\Ga(\frac{a}{N})\,\Ga(\frac{b}{N})}{\Ga(\frac{a+b}{N})}.
\end{equation}
Образ период-функционала $\gamma\mapsto\int_\gamma\eta$ на
$H_1(C_N,\ZZ)$ содержится в решётке
$\frac{\Omega_{a,b}}{N}\,\ZZ[\zeta_N]$. В частности, $\Omega_{a,b}>0$
вещественно.
\end{lemma}

\begin{proof}
Тождество~\eqref{eq:beta}: подстановка $t=x^N$ даёт $dt=Nx^{N-1}dx$,
и потому $x^i\,dx=\tfrac1N t^{\frac{i+1}{N}-1}dt$; далее
$y^j/y^{N-1}=y^{j-N+1}=(1-t)^{\frac{j+1}{N}-1}$, так как $y^N=1-t$.
Собирая, получаем~\eqref{eq:beta} с $a=i+1$, $b=j+1$. Интегрирование
по $\gamma$ есть интегрирование формы~\eqref{eq:beta} по пути
$\delta$ с выбранными ветвями степеней; каждая полная петля вокруг
$t=0$ умножает ветвь $t^{1/N}$ на $\zeta_N$ и, значит, интеграл от
$t^{a/N-1}dt$ --- на $\zeta_N^{a}$; аналогично петля вокруг $t=1$
умножает $(1-t)^{1/N}$ на $\zeta_N$ и интеграл --- на $\zeta_N^{b}$
(обход $t$ вокруг $1$ есть обход $1-t$ вокруг $0$ с сохранением
ориентации). Путь из опорной точки в опорную с данными обхода
$(r,s)$ даёт, следовательно, множитель $\zeta_N^{ra+sb}$ на
эталонном интеграле
\[
  \frac1N\int_0^1 t^{\frac{a}{N}-1}(1-t)^{\frac{b}{N}-1}\,dt
  \;=\;\frac1N\,\Bfun\!\Bigl(\frac{a}{N},\frac{b}{N}\Bigr),
\]
где эталонный путь --- отрезок $(0,1)$ с главными ветвями. Это и
есть~\eqref{eq:closedform}. Положительность $\Omega_{a,b}>0$ ---
стандартное свойство бета-интеграла при $a/N>0$, $b/N>0$. Содержание
образа период-функционала в решётке $\Omega_{a,b}\ZZ[\zeta_N]/N$
следует из того, что любой цикл есть целочисленная комбинация
замкнутых цепей описанного вида (подъёмы рёбер графа $\Gamma_N$
порождают $H_1$ после замыкания --- леммы~\ref{lem:gridgraph}
и~\ref{lem:capping}), а интеграл по каждому ребру есть разность двух
значений вида~\eqref{eq:closedform}.
\end{proof}

\begin{corollary}[$\mu_N\times\mu_N$-эквивариантность]\label{cor:equivar}
Перенос цикла действием $(\zeta_N^u,\zeta_N^v)\in\mu_N\times\mu_N$
умножает период собственной формы $(a,b)$ на $\zeta_N^{ua+vb}$:
\[
  \int_{\gamma\cdot(\zeta^u,\zeta^v)}\eta_{a,b}
  \;=\;\zeta_N^{\,ua+vb}\;\int_\gamma\eta_{a,b}.
\]
\end{corollary}

\begin{proof}
Действие поднимает координату $t=x^N$ тождественно, поэтому отображает
пути с данными обхода $(r,s)$ в пути с данными обхода
$(r+u,\,s+v)$ (ветви $t^{1/N}$ и $(1-t)^{1/N}$ умножаются на
$\zeta^u$ и $\zeta^v$). Остаётся применить замкнутую
форму~\eqref{eq:closedform}: фаза
$\zeta^{(r+u)a+(s+v)b}=\zeta^{ua+vb}\zeta^{ra+sb}$.
\end{proof}

\begin{remark}
Замкнутая форма~\eqref{eq:closedform} выполняет обе функции, ради
которых в раннем плане стенда предусматривались L-значения и PSLQ:
она фиксирует фазу \emph{явно} (никакая «общая фаза» не ищется) и
заменяет численный поиск соотношений \emph{тождественным выводом} из
$\Ga$-тождеств (раздел~\ref{sec:gammaid}). Вычислительная сверка
замкнутой формы с независимым прямым интегрированием (tanh-sinh по
гладким кускам, без $\Ga$-функций) --- проверка V2 протокола;
максимальное относительное отклонение на выборке $69$ тестов для
$N=15,30$ составляет $3.9\cdot10^{-36}$ (раздел~\ref{sec:protocol}).
\end{remark}

\section{$\Ga$-тождества (лемма~2)}\label{sec:gammaid}

Здесь собирается весь аналитический инструментарий стенда --- два
классических функциональных тождества гамма-функции. Принципиально,
что больше ничего не требуется: все соотношения между используемыми
$\Ga$-произведениями выводятся из этих двух тождеств, что и делает
изложение самодостаточным.

\begin{lemma}[формула отражения]\label{lem:reflection}
Для всех $z\in\CC\setminus\ZZ$
\begin{equation}\label{eq:reflection}
  \Ga(z)\,\Ga(1-z)\;=\;\frac{\pi}{\sin \pi z}.
\end{equation}
\end{lemma}

\begin{proof}[Доказательство (полный вывод --- приложение~\ref{app:gamma})]
Ключевое звено --- интеграл Эйлера
\[
  \int_0^\infty \frac{u^{z-1}}{1+u}\,du \;=\; \frac{\pi}{\sin\pi z}
  \qquad (0<\re z<1),
\]
доказываемый вычетами по ключевому контуру. Замена
$t=\frac{u}{1+u}$ приводит этот интеграл к
$\Bfun(z,1-z)=\Ga(z)\Ga(1-z)$. Сопоставление даёт~\eqref{eq:reflection}.
\end{proof}

\begin{lemma}[формула умножения Гаусса]\label{lem:gauss}
Для целых $m\ge2$ и $z$ с $mz\notin\ZZ_{\le0}$
\begin{equation}\label{eq:gauss}
  \Ga(z)\,\Ga\!\Bigl(z+\frac1m\Bigr)\cdots
  \Ga\!\Bigl(z+\frac{m-1}{m}\Bigr)
  \;=\;(2\pi)^{\frac{m-1}{2}}\;m^{\frac12-mz}\;\Ga(mz).
\end{equation}
\end{lemma}

\begin{proof}[Доказательство (полный вывод --- приложение~\ref{app:gamma})]
Аргумент через рекуррентность и формулу Стирлинга: отношение левой и
правой частей есть $1$-периодическая голоморфная функция, стремящаяся
к $1$ при $\re z\to+\infty$ в вертикальных полосах конечной ширины;
периодичность переносит предел на всю полосу, откуда тождество.
\end{proof}

\begin{lemma}[инвентарь соотношений]\label{lem:relations}
Любое соотношение между элементами семейства
$\{\Omega_{a,b}=\Bfun(a/N,b/N)\}_{(a,b)\in T_N}$, используемое в
настоящем тексте, есть следствие тождеств~\eqref{eq:reflection}
и~\eqref{eq:gauss}. В частности:
\begin{enumerate}[leftmargin=2em, itemsep=0.15em]
\item если $a+b=N$, то
  $\Omega_{a,b}=\Ga(a/N)\Ga(b/N)=\pi/\sin(\pi a/N)$ --- чистая дробная
  доля $\pi$;
\item если $a+b>N$, то
  $\Omega_{a,b}=\dfrac{N}{a+b-N}\cdot
  \dfrac{\Ga(a/N)\,\Ga(b/N)}{\Ga\bigl(\frac{a+b-N}{N}\bigr)}$
  --- редукция к тройке аргументов с суммой $<N$;
\item если $a+b<N$, то $\Omega_{a,b}$ есть $\Bfun$-значение в
  «канонической» позиции и редукции не требует.
\end{enumerate}
\end{lemma}

\begin{proof}
Пункт~1 --- прямое применение~\eqref{eq:reflection} к
$b/N=1-a/N$. Пункт~2 --- рекуррентность
$\Ga\bigl(\frac{a+b}{N}\bigr)=\Ga\bigl(1+\frac{a+b-N}{N}\bigr)
=\frac{a+b-N}{N}\,\Ga\bigl(\frac{a+b-N}{N}\bigr)$. Пункт~3 --- по
определению. Поскольку никаких иных преобразований значений
$\Ga(k/N)$ в изложении не производится, всякий вывод о соотношениях
сводится к последовательному применению этих трёх правил и тождеств
\eqref{eq:reflection}--\eqref{eq:gauss}.
\end{proof}

\begin{statusbox}[границы утверждаемого]
Лемма~\ref{lem:relations} утверждает ровно то, что нужно для стенда:
все \emph{используемые} соотношения выводятся из отражения и
умножения. Проверка достаточности порождающего набора выполняется не
поиском соотношений, а прямой проверкой ранга явно построенной матрицы
функционалов (V7 протокола: численный ранг равен $g$ --- $91$ и
$406$ --- при хорошем числе обусловленности $8.6$ и $18.2$).
\end{statusbox}
